\documentclass[dvipsnames]{beamer}
\usepackage[utf8]{inputenc}      % UTF-8编码支持
\usepackage[T1]{fontenc}         % 字体编码
\usepackage{amsmath,amssymb}  % 数学符号和环境（beamer自带proof/\qed，无需amsthm，避免冲突）
\usepackage{graphicx}            % 图片插入
\usepackage{xcolor}              % 颜色支持
\usepackage{tikz}                % TikZ绘图
\usepackage{framed}              % 框架环境
\usepackage{tcolorbox}           % 彩色盒子（用于定理环境）
\usepackage{verbatim}
\tcbuselibrary{skins, breakable, theorems}   % tcolorbox的定理/皮肤库（enhanced、drop shadow 需要 skins）
\renewcommand{\familydefault}{\sfdefault}
\usetheme{Madrid}
\usecolortheme{dolphin}
\makeatletter
\setbeamertemplate{footline}{%
  \leavevmode%
  \hbox{%
    \begin{beamercolorbox}[wd=.37\paperwidth,ht=2.25ex,dp=1ex,center]{author in head/foot}%
      \usebeamerfont{author in head/foot}\insertshortauthor\expandafter\ifblank\expandafter{\beamer@shortinstitute}{}{~~(\insertshortinstitute)}
    \end{beamercolorbox}%
    \begin{beamercolorbox}[wd=.37\paperwidth,ht=2.25ex,dp=1ex,center]{title in head/foot}%
      \usebeamerfont{title in head/foot}\insertshorttitle
    \end{beamercolorbox}%
    \begin{beamercolorbox}[wd=.26\paperwidth,ht=2.25ex,dp=1ex,center]{date in head/foot}%
      \usebeamerfont{date in head/foot}\usebeamertemplate{page number in head/foot}\hspace*{2ex}%
    \end{beamercolorbox}}%
  \vskip0pt%
}
\makeatother
\setbeamercolor{titlelike}{bg=Emerald!20, fg=JungleGreen!75!NavyBlue}
\setbeamerfont{title}{series=\bfseries}
%title page
\setbeamercolor{subtitle}{fg=Emerald!70!white}
\setbeamercolor{author}{fg=JungleGreen!80!NavyBlue}
\setbeamercolor{institute}{fg=JungleGreen!80!NavyBlue}
\setbeamercolor{date}{fg=Emerald!70!white}
% footline
\setbeamercolor{footline}{fg=JungleGreen!80!NavyBlue}
\setbeamercolor{author in head/foot}{bg=Emerald!10, fg=JungleGreen!80!NavyBlue}
\setbeamercolor{title in head/foot}{bg=Emerald!18, fg=JungleGreen!80!NavyBlue}
\setbeamercolor{page number in head/foot}{bg=Emerald!25, fg=JungleGreen!60!NavyBlue}
\setbeamercolor{date in head/foot}{fg=Emerald, bg=Emerald!25}
\setbeamercolor{navigation symbols}{fg=NavyBlue!50}
% table of contents / items
\setbeamercolor{structure}{fg=JungleGreen!80!NavyBlue}
\setbeamertemplate{section in toc}[sections numbered]
\setbeamercolor{section in toc}{fg=JungleGreen!80!NavyBlue}
\setbeamercolor{section number projected}{bg=BlueGreen, fg=NavyBlue}
\setbeamercolor{item projected}{bg=BlueGreen, fg=NavyBlue}
\setbeamercolor{item}{fg=JungleGreen!80!NavyBlue}
\setbeamercolor{subitem}{fg=JungleGreen!80!NavyBlue}
\setbeamercolor{subsubitem}{fg=JungleGreen!80!NavyBlue}
\AtBeginDocument{
  \color{JungleGreen!80!NavyBlue}
}

% Theorem-like environments not predefined by beamer
\theoremstyle{plain}
\newtheorem{proposition}[theorem]{\translate{Proposition}}
\theoremstyle{definition}
\newtheorem{exercise}[theorem]{\translate{Exercise}}
\newtheorem{remark}[theorem]{\translate{Remark}}

\title{MATH1860J Recitation Class 01}
\subtitle{Foundations: Chapters 01-05}
\author{Shuhan Yu}
\institute{SJTU Global College}
\date{Fall Term 2026}

\begin{document}

\AtBeginSection[]
{
  {
    \setbeamercolor{background canvas}{bg=Aquamarine!75}
    \setbeamercolor{normal text}{fg=white}
    \setbeamercolor{frametitle}{fg=white}

    \begin{frame}[plain]
      \vfill
      {\bfseries\Large\color{white!60}
        \thesection.~\insertsection

        \vspace{0.3em}
        \noindent
        \textcolor{white}{\rule{0.58\paperwidth}{1pt}}%
        \textcolor{white!45}{\rule{0.27\paperwidth}{1pt}}
      }
      \vfill
    \end{frame}
  }
}

\begin{frame}
    \titlepage
\end{frame}

\begin{frame}{Outline}
    \tableofcontents
\end{frame}

\section{Introduction}
\begin{frame}{Introduction}
About me:
\begin{itemize}
    \item \textbf{Shuhan Yu}: Sophomore student majoring in ECE at GC.
    \item \textbf{Email:} \texttt{barca0927@sjtu.edu.cn}
    \item \textbf{Office Hours:} 
    \item \textbf{Interest:} Quantitative Trading, Machine Learning, and Geoscience.
\end{itemize}
\vspace{0.5em}
About Honors Mathematics Sequence (8 series):
\begin{itemize}
    \item Emphasize on rigorous mathematical reasoning and proof.
    \item Write every step of your solution clearly and logically.
    \item Don't be frightened by ``Honor'' in the course title. Remember that 5 series is also an honor series.
    \item Sample exams are of great help. Go through them before exams.
\end{itemize}
Feel free to interrupt and ask questions if you have any doubt during RCs!
\end{frame}

\section{Logic}
\begin{frame}{Propositional Logic}

\begin{definition}
A \textbf{statement} (also called a \textbf{proposition}) is a declarative sentence 
that can be unambiguously classified as either \emph{true} or \emph{false}.
\end{definition}

\begin{itemize}
  \item \textbf{Logical operations:}
    \begin{itemize}
      \item Negation (unary operation): $\lnot A$
      \item Conjunction: $A \land B$
      \item Disjunction: $A \lor B$
      \item Implication: $A \to B$
      \item Converse: $B \to A$
      \item Equivalence: $A \leftrightarrow B$
    \end{itemize}
    \textcolor{Aquamarine}{Binary operations return a compound statement.}
  \item \textbf{Special propositions:}
    \begin{itemize}
      \item Tautology: always true
      \item Contradiction: always false
      \item Vacuous truth.
    \end{itemize}
\end{itemize}

\end{frame}

\begin{frame}{Propositional Logic}

The definition of logical operations is often given by a \textbf{truth table}.  Consider the implication $A \to B$.  
It can be rewritten in several equivalent forms (where $\equiv$ denotes logical equivalence):
\[
A \to B \;\;\equiv\;\; \lnot A \lor B 
\;\;\equiv\;\; \lnot(A \land \lnot B) 
\;\;\equiv\;\; \lnot B \to \lnot A.
\]

\begin{table}[htbp]
\centering
\renewcommand{\arraystretch}{1.3}
\begin{tabular}{c c | c c c c}
\hline
$A$ & $B$ & $A \to B$ & $\lnot A \lor B$ & $\lnot(A \land \lnot B)$ & $\lnot B \to \lnot A$ \\
\hline
T & T & T & T & T & T \\
T & F & F & F & F & F \\
F & T & T & T & T & T \\
F & F & T & T & T & T \\
\hline
\end{tabular}
\end{table}
With the basic operation laws established, we can use them to simplify and analyze more complicated propositions.
\end{frame}

\begin{frame}{Propositional Logic}
\textbf{Remark:} (Important for Proof)
\begin{itemize}
    \item $A \to B \equiv \lnot B \to \lnot A$: equivalence with the \textbf{contrapositive}.  
    \item $A \to B \equiv \lnot(A \land \lnot B)$: useful in \textbf{proof by contradiction}.  
    \item $A \to B \equiv \lnot A \lor B$: converts the implication into a disjunction, often used in algebraic manipulation.  
\end{itemize}
\vspace{0.5em}
\textbf{Important logical equivalences and laws:}
\begin{itemize}
  %\item \textbf{Identity Laws:} $p \land \text{T} \equiv p, \; p \lor \text{F} \equiv p$
  %\item \textbf{Domination Laws:} $p \lor \text{T} \equiv \text{T}, \; p \land \text{F} \equiv \text{F}$
  %\item \textbf{Double Negation:} $\lnot(\lnot p) \equiv p$
  \item \textbf{Commutativity:} $p \lor q \equiv q \lor p, \; p \land q \equiv q \land p$
  \item \textbf{Associativity:} $(p \lor q) \lor r \equiv p \lor (q \lor r)$
  \item \textbf{Distributivity:} 
  \begin{align*}
    p \lor (q \land r) &\equiv (p \lor q) \land (p \lor r)\\
    p \land (q \lor r) &\equiv (p \land q) \lor (p \land r)
  \end{align*}
  \item \textbf{De Morgan's Laws:} 
  \begin{align*}
    \lnot(p \land q) &\equiv \lnot p \lor \lnot q\\
    \lnot(p \lor q) &\equiv \lnot p \land \lnot q
  \end{align*}
\end{itemize}

\end{frame}

\begin{frame}{Predicate Logic}

A \textbf{statement} can be expressed as:
\[
\textbf{statement} = \textbf{quantifier} + \textbf{predicate (statement frame)}.
\]

\begin{definition}
A \textbf{predicate} is an expression $E(x)$ that becomes a statement when 
$x$ is replaced by an object from \textbf{a specified class} of objects.
\vspace{0.5em}
\begin{itemize}
    \item \textbf{Universal quantifier} $\forall$: ``for all''.  
    The expression $\forall x \in X: E(x)$ means ``Every object $x$ in $X$ satisfies property $E$''.
    \vspace{0.5em}
    \item \textbf{Existential quantifier} $\exists$: ``there exists''.  
    The expression $\exists x \in X: E(x)$ means ``There exists at least one object $x$ in $X$ with property $E$''.
    \vspace{0.5em}
    \item \textbf{Unique existential} $\exists!$: exactly one such object exists.
\end{itemize}
\end{definition}

\end{frame}


\begin{frame}{Predicate Logic}

In all cases, \textbf{the order of the quantifiers is significant.}
Thus, 
$$
\forall x \ \exists y:\, E(x,y) \quad \nLeftrightarrow  \quad \exists y \ \forall x:\, E(x,y)
$$
\begin{itemize}
    \item Former, for all $x$ there exists some $y$ such that $E(x,y)$ is true. 
    \item Latter, there exists a fixed $y$ such that $E(x,y)$ is true for all $x$. 
\end{itemize}
\vspace{0.5em}
\end{frame}

\begin{frame}[t]
\frametitle{Exercise}

\begin{block}{Exercise 01}
    \begin{enumerate}
        \item Show that \[(A \land B) \to C \;\;\equiv\;\; A \to B \to C \;\;\equiv\;\; B \to A \to C.\]
        \item Give the negations of the statements: $$ \forall x \in X, \exists y \in Y:\, E(x,y,z)$$
        $$\forall \epsilon > 0, \exists \delta > 0:\, |x - a| < \delta \implies |f(x) - f(a)| < \epsilon$$
    \end{enumerate}
\end{block}
\textbf{Remarks:}\\
Using the quantifiers $\exists$ and $\forall$, negation becomes a purely “mechanical” process in which the symbols $\exists$ and $\forall$ (as well as $\wedge$ and $\vee$) are interchanged (without changing the order), and statements are negated.
\end{frame}

\section{Set Theory}

\begin{frame}{Set Basics and Operations}
    Sets are collections of objects, called \textbf{elements} or \textbf{members}.
    \begin{itemize}
    \item Define by Listing: \textit{A} = \{1, 2, 3, 4\}
    \item Define by Predicate: \textit{A} = \{$x \in \mathbb{N} \mid 1 \leq x \leq 4$\}
    \item Set $A = B$ iff $A\subseteq B$ and $B \subseteq A$. (Useful in Proof)
    \end{itemize}
\vspace{0.5em}
Let \(A = \{x \mid P_1(x)\}\) and \(B = \{x \mid P_2(x)\}\).  
\begin{itemize}
\item \text{Union:} $A \cup B := \{x \mid P_1(x) \lor P_2(x)\}$
\item \text{Intersection:} $A \cap B := \{x \mid P_1(x) \land P_2(x)\}$
\item \text{Difference:} $A \setminus B := \{x \mid P_1(x) \land \lnot P_2(x)\}$.
\item \text{Complement:} $A^c (\overline{A}) := \mathbb{U} \setminus A := \{x \in \mathbb{U} \mid \lnot P_1(x)\}$.
\end{itemize}
If \(A \cap B = \emptyset\), we say that \(A\) and \(B\) are \textbf{disjoint}.  \\
\vspace{0.5em}
\textbf{Remark:} Distinguish between \textbf{Difference}(\(A \setminus B\)) and \textbf{Quotient}(\(A / B\))    
\end{frame}


\begin{frame}{Set Basics and Operations}

\begin{block}{Exercise 02}
Verify the following set-theoretic identities:

\begin{enumerate}
    \item $(A \cup B)^c = A^c \cap B^c$
    \item $(A \cap B)^c = A^c \cup B^c$
    \item $A \cap (B \cup C) = (A \cap B) \cup (A \cap C)$
    \item $A \cup (B \cap C) = (A \cup B) \cap (A \cup C)$
    \item $(A \cup B) \setminus C = (A \setminus C) \cup (B \setminus C)$
    \item $(A \cap B) \setminus C = (A \setminus C) \cap (B \setminus C)$
    \item $A \setminus (B \cup C) = (A \setminus B) \cap (A \setminus C)$
    \item $A \setminus (B \cap C) = (A \setminus B) \cup (A \setminus C)$
\end{enumerate}

\textbf{Hint:} Use element-wise reasoning and draw a Venn diagram.
\end{block}

\end{frame}


\begin{frame}{Power Set and Cardinality}

\begin{itemize}
    \item \textbf{Power Set:} $P(A) = \{ X \mid X \subseteq A \}$
    \item \textbf{Cardinality:} $|A|$ is the number of elements in set $A$
    \item \textbf{Finite Set:} A set with a finite number of elements
    \item \textbf{Infinite Set:} A set with an infinite number of elements
    \item \textbf{Countable Set:} A set that can be put into one-to-one correspondence with the natural numbers
    \item \textbf{Uncountable Set:} A set that cannot be put into one-to-one correspondence with the natural numbers
    \item \textbf{Cardinality of Power Set:} $|P(A)| = 2^{|A|}$ if $A$ is finite
\end{itemize}

\end{frame}

\begin{frame}{Ordered Pairs and Cartesian Product}
\begin{itemize}
    \item \textbf{Ordered Pair:} $(a,b)$
    \item \textbf{Cartesian Product:} $A \times B = \{ (a,b) \mid a \in A, b \in B \}$
    \item \textbf{n-tuple}: generalization of ordered pair, $(a_1, a_2, \ldots, a_n)$
    \item Construction of ordered pair and natural numbers: $(a,b) = \{\{a\}, \{a,b\}\}; 0 = \{\emptyset\}, 1 = \{\emptyset\}, 2 = \{\emptyset, \{\emptyset\}\}, ...$
\end{itemize}
\vspace{0.5em}
- Note that \((a,b) = (c,d) \iff a = c \wedge b = d\).\\
- In general, $A \times B \neq B \times A$.
\end{frame}

\begin{comment}
\begin{frame}{Relation}

\begin{definition}
A \textbf{binary relation} on a set $X$ is a subset $R \subseteq X \times X$.  
Instead of writing $(x,y) \in R$, we usually write $xRy$ or $x \sim_R y$.
\end{definition}

\textbf{Properties of Relations:}  
\begin{itemize}
    \item \textbf{Reflexive:} $R$ is reflexive if $xRx$ for all $x \in X$.
    \item \textbf{Transitive:} $R$ is transitive if $(xRy) \wedge (yRz) \implies xRz$.
    \item \textbf{Symmetric:} $R$ is symmetric if $xRy \implies yRx$.
\end{itemize}

Let $Y \subseteq X$ be nonempty. The \textbf{restriction} of $R$ to $Y$ is  
\[
R_Y := (Y \times Y) \cap R.
\]  
Clearly, $xRYy$ if and only if $x, y \in Y$ and $xRy$.  
Usually, we simply write $R$ instead of $R_Y$ when the context is clear.

\end{frame}

\begin{frame}{Equivalence Relation}
\begin{definition}[Equivalence Relation]
A relation $\sim$ on a set $X$ is called an \textbf{equivalence relation} if it is:
\begin{itemize}
    \item \textbf{Reflexive:} $x \sim x$ for all $x \in X$,
    \item \textbf{Symmetric:} $x \sim y \implies y \sim x$,
    \item \textbf{Transitive:} $(x \sim y \wedge y \sim z) \implies x \sim z$.
\end{itemize}
\end{definition}

\begin{definition}[Equivalence Class and Quotient Set]
For $x \in X$, the \textbf{equivalence class} of $x$ is
\[
[x] := \{ y \in X \mid y \sim x \}.
\]  
The set of all equivalence classes of $X$ is
\[
X / \sim := \{ [x] \mid x \in X \},
\]
called \textbf{$X$ modulo $\sim$}.
\end{definition}
\end{frame}

\begin{frame}{Partition}
\begin{definition}[Partition]
A \textbf{partition} of a set $X$ is a subset $\mathcal{A} \subseteq \mathcal{P}(X) \setminus \{\emptyset\}$ such that:
\begin{itemize}
    \item For each $x \in X$, there is a unique $A \in \mathcal{A}$ with $x \in A$,
    \item The elements of $\mathcal{A}$ are pairwise disjoint, and $\bigcup \mathcal{A} = X$.
\end{itemize}
\end{definition}

\begin{proposition}
Let $\sim$ be an equivalence relation on $X$. Then $X / \sim$ is a partition of $X$.
\end{proposition}

For all $x \in X$, $x \in [x]$, so $X = \bigcup_{x \in X} [x]$.  
If $z \in [x] \cap [y]$, then $z \sim x$ and $z \sim y$, hence $x \sim y$, which implies $[x] = [y]$.  
Thus, equivalence classes are either identical or disjoint.
\end{frame}

\begin{frame}{Order Relations}
\begin{definition}[Partial Order]
A relation $\le$ on a set $X$ is a \textbf{partial order} if it is:
\begin{itemize}
    \item \textbf{Reflexive:} $x \le x$ for all $x \in X$,
    \item \textbf{Transitive:} $x \le y \wedge y \le z \implies x \le z$,
    \item \textbf{Antisymmetric:} $x \le y \wedge y \le x \implies x = y$.
\end{itemize}
The pair $(X, \le)$ is called a \textbf{partially ordered set}. If the order is clear from context, we simply write $X$.
\end{definition}

\begin{definition}[Total Order]
If, in addition, 
\[
\forall x, y \in X: (x \le y) \vee (y \le x),
\] 
then $\le$ is a \textbf{total order}, and $(X, \le)$ is called a \textbf{totally ordered set}.
\end{definition}
\end{frame}

\begin{frame}{Order Relations}

\textbf{Useful notations:}\quad
$x \ge y \;:\iff\; y \le x$,\quad
$x < y \;:\iff\; x \le y \wedge x \neq y$,\quad
$x > y \;:\iff\; y < x$.

\begin{remark}
If \(X\) is \textbf{totally ordered}, then for each \(x, y \in X\), exactly one of
\(x < y\), \(x = y\), \(x > y\) holds.
\end{remark}

\begin{remark}
If \(X\) is \textbf{partially ordered} but not totally ordered, there exist elements \(x, y \in X\) which are \textbf{incomparable}, i.e., neither \(x \le y\) nor \(y \le x\) holds.
\end{remark}

\end{frame}
\end{comment}

\begin{frame}{Naive Set Theory and Russell's Paradox}
\textbf{Naive Set Theory:} You can define a set by any property. A set is any collection of distinct objects.
\textbf{Problem: Russell's Paradox} arises when considering the set of all sets that do not contain themselves:
\[
S = \{ x \mid x \text{ is a set and } x \notin x \}
\]

\textbf{Question:} Is $S \in S$?

\begin{itemize}
    \item If $S \in S$, then by definition $S \notin S$.
    \item If $S \notin S$, then by definition $S \in S$.
\end{itemize}

\textbf{Conclusion:} Contradiction arises. Not every property can define a set freely. $\Rightarrow$ \textbf{Implication:} Motivates axiomatic set theory (e.g., Zermelo-Fraenkel Set Theory, ZFC) to avoid such paradoxes.
\end{frame}

\begin{frame}{Cantor's Paradox (Supplement)}
\textbf{Prove that ``the set of all sets'' does not exist.} \\
\vspace{0.5em}
\textbf{Cantor's Theorem:} For any set $A$,
\[
|A| < |P(A)|,
\]
where $P(A)$ is the power set of $A$, i.e., the set of all subsets of $A$.

\textbf{Consider:} The ``set of all sets'' $U$.

\begin{itemize}
    \item By definition, $U$ contains every set, so $P(U) \subseteq U$.
    \item By Cantor's Theorem, $|U| < |P(U)|$.
\end{itemize}

\textbf{Contradiction:} $|U| \ge |P(U)|$ from definition, but $|U| < |P(U)|$ by Cantor.  

\textbf{Conclusion:} There is no ``set of all sets''. This prevents the paradox.
\end{frame}

\begin{frame}{Zermelo-Fraenkel Set Theory (ZF)}
\footnotesize
\begin{enumerate}
    \item \textbf{Axiom of Extensionality:} Two sets are equal if they have the same elements.
    \item \textbf{Axiom of Empty Set:} There exists a set with no elements, denoted $\emptyset$.
    \item \textbf{Axiom of Pairing:} For any sets $a$ and $b$, there exists a set $\{a,b\}$.
    \item \textbf{Axiom of Union:} For any set of sets $A$, there exists a set $\bigcup A$ containing all elements of elements of $A$.
    \item \textbf{Axiom of Power Set:} For any set $A$, there exists a set $P(A)$ of all subsets of $A$.
    \item \textbf{Axiom of Infinity:} There exists a set containing $\emptyset$ and closed under successor.
    \item \textbf{Axiom of Replacement:} The image of a set under any definable function is also a set.
    \item \textbf{Axiom of Separation (Specification):} For any set $A$ and property $P$, $\{x \in A \mid P(x)\}$ is a set.
    \item \textbf{Axiom of Regularity (Foundation):} Every non-empty set $A$ contains an element disjoint from $A$.
\end{enumerate}
\end{frame}

\begin{frame}{Exercise}
\begin{block}{Exercise 03}
    Let $A,B,C$ be three statements. Use truth table to prove that
    $$
    (A \lor B) \land C \equiv (A \land C) \lor (B \land C)
    $$
    Let $A,B,C$ be three sets. Prove that (may use Venn Diagram)
    $$
    (A \cup B) \cap C = (A \cap C) \cup (B \cap C)
    $$
\end{block}
\end{frame}


\section{Natural Numbers and Mathematical Induction}
\begin{comment}
\begin{frame}{Algebra Structure}
\begin{definition}[Internal and External Composition Laws]
Let $S$ and $S'$ be two sets.
\begin{enumerate}
    \item An \textbf{internal composition law} $\circ$ is a map
    \[
        S \times S \to S, \quad (x,y) \mapsto x \circ y.
    \]
    \item An \textbf{external composition law} $\ast$ is a map
    \[
        S' \times S \to S, \quad (\alpha,x) \mapsto \alpha \ast x.
    \]
\end{enumerate}
\end{definition}
\end{frame}

\begin{frame}{Algebra Structure}
\begin{definition}[Group]
A \textbf{group} is a pair $(G,+)$ consisting of a set $G$ and an internal composition law $+$ such that:
\begin{enumerate}
    \item \textbf{Associativity:} $a + (b + c) = (a + b) + c$ for all $a,b,c \in G$.
    \item \textbf{Existence of identity:} there exists $e_+ \in G$ such that $a + e_+ = e_+ + a = a$ for all $a \in G$.
    \item \textbf{Existence of inverse:} for each $a \in G$ there exists $a^{-1} \in G$ such that $a + a^{-1} = a^{-1} + a = e_+$.
\end{enumerate}
A group is \textbf{abelian} if it also satisfies:
\begin{enumerate}
    \item[iv] \textbf{Commutativity:} $a + b = b + a$ for all $a,b \in G$.
\end{enumerate}
\end{definition}
\end{frame}

\begin{frame}{Algebra Structure}
\begin{definition}[Ring]
A \textbf{ring} is a triple $(R,+,\times)$ consisting of a set $R$ and two internal composition laws $+$ and $\times$, such that:
\begin{enumerate}
    \item $(R,+)$ is an abelian group.
    \item \textbf{Associativity:} $a \times (b \times c) = (a \times b) \times c$ for all $a,b,c \in R$.
    \item \textbf{Multiplicative identity:} there exists $e_\times \in R$ such that $a \times e_\times = e_\times \times a = a$ for all $a \in R$.
    \item \textbf{Distributivity:} $a \times (b + c) = (a \times b) + (a \times c)$ and $(b + c) \times a = (b \times a) + (c \times a)$ for all $a,b,c \in R$.
\end{enumerate}
A ring is \textbf{commutative} if it also satisfies:
\begin{enumerate}
    \item[v] $a \times b = b \times a$ for all $a,b \in R$.
\end{enumerate}
\end{definition}
\end{frame}

\begin{frame}{Algebra Structure}
\begin{definition}[Field]
A \textbf{field} is a commutative ring $(F,+,\times)$ with additive identity $e_+$ and multiplicative identity $e_\times$ such that:
\begin{enumerate}
    \item $e_+ \neq e_\times$.
    \item For every $a \in F \setminus \{ e_+ \}$, there exists $a^{-1} \in F$ such that $a \times a^{-1} = e_\times$.
\end{enumerate}
The \textbf{characteristic} of a field is the smallest $n \in \mathbb{N}^*$ such that $n \cdot 1 = 0$, if it exists.
\end{definition}
\end{frame}
\end{comment}

\begin{frame}{Natural Numbers}
\textbf{Natural Numbers from ZF Set Theory}

\begin{itemize}
    \item By the \textbf{Axiom of Infinity}, there exists a set $I$ such that:
    \begin{itemize}
        \item $\emptyset \in I$
        \item For each $x \in I$, $S(x) := x \cup \{x\} \in I$
    \end{itemize}
    \item \textbf{Define} natural numbers \textbf{recursively}:
    \[
        0 := \emptyset, \quad 1 := S(0) = \{\emptyset\}, \quad 2 := S(1) = \{\emptyset, \{\emptyset\}\}, \dots
    \]
    \item The set of natural numbers $\mathbb{N}$ is the smallest set containing $0$ and closed under $S$.
    \item Ordering: The natural numbers are well-ordered.
\end{itemize}
\vspace{0.5em}
\textbf{Comment:}
$\{\cdot\}$ is packaging and $\cup$ is stacking. \\
Each new number encompasses all its previous history and incorporates ``the previous self'' as a new element. This is why natural numbers have sequentiality.
\end{frame}

%--------------------------------
\begin{frame}{Natural Numbers}
\textbf{Recursive Definition of Addition on $\mathbb{N}$}

\begin{itemize}
    \item \textbf{Base case:} 
    \[
        \forall n \in \mathbb{N}, \quad n + 0 := n
    \]
    \item \textbf{Successor step:} 
    \[
        \forall n,m \in \mathbb{N}, \quad n + S(m) := S(n + m)
    \]
\end{itemize}
\textit{Explanation:} $S(m) = m \cup \{m\}$ is the successor of $m$, so addition recursively “adds 1” using the successor operation. \\
\vspace{0.5em}
\textbf{Properties of addition:}
\begin{itemize}
    \item Commutativity: $n + m = m + n$
    \item Associativity: $(n + m) + k = n + (m + k)$
    \item Identity/Existence of a neutral element: $n + 0 = n$     
\end{itemize}
\end{frame}

%--------------------------------
\begin{frame}{Natural Numbers}
\textbf{Recursive Definition of Multiplication on $\mathbb{N}$}

\begin{itemize}
    \item \textbf{Base case:} 
    \[
        \forall n \in \mathbb{N}, \quad n \cdot 0 := 0
    \]
    \item \textbf{Successor step:} 
    \[
        \forall n,m \in \mathbb{N}, \quad n \cdot S(m) := (n \cdot m) + n
    \]
\end{itemize}
\textit{Explanation:} Multiplication is defined as repeated addition, consistent with the intuitive meaning of natural numbers. \\
\vspace{0.5em}
\textbf{Properties of multiplication:}
\begin{itemize}
    \item Commutativity: $n \cdot m = m \cdot n$
    \item Associativity: $(n \cdot m) \cdot k = n \cdot (m \cdot k)$
    \item Identity/Existence of a neutral element: $n \cdot 1 = n$  
    \item Distributivity: $n \cdot (m + k) = (n \cdot m) + (n \cdot k)$   
\end{itemize}
\end{frame}

\begin{frame}{Mathematical Induction}

\begin{definition}[Principle of Mathematical Induction]
Let $A(n)$ be a statement depending on $n \in \mathbb{N}$ and let $n_0 \in \mathbb{N}$.  
To prove $A(n)$ for all $n \ge n_0$, it suffices to verify:

\begin{enumerate}
    \item \textbf{Base Case:} $A(n_0)$ is true.
    \item \textbf{Inductive Step:} For all $n \ge n_0$, if $A(n)$ is true, then $A(n+1)$ is true:
    \[
    \forall n \in \mathbb{N},\; n \ge n_0 : \quad A(n) \implies A(n+1).
    \]
\end{enumerate}

\textbf{Conclusion:} If both conditions hold, $A(n)$ is true for all $n \ge n_0$.
\end{definition}
\textbf{Remark:} The principle of induction is essentially part of the definition of the natural numbers $\mathbb{N}$.
\end{frame}

\begin{frame}{Mathematical Induction}

\begin{block}{Exercise 04}
Suppose $M$ is a finite set of men.

\textbf{Proposition:} For any $a,b \in M$, $a = b$.

\textbf{Proof (by induction on $|M|$):}
\begin{enumerate}
    \item \textbf{Base Case:} If $|M|=1$, the claim is trivially true.
    \item \textbf{Inductive Step:} Assume true for $n$ men.  
    Let $|M| = n+1$ and $a,b \in M$.  
    Define $M_a = M \setminus \{a\}$, $M_b = M \setminus \{b\}$.  
    Take $c \in M_a \cap M_b$. Then $a = c$, $b = c \implies a = b$.
\end{enumerate}

\textbf{Question:} Where is the flaw in this proof?
\end{block}
\begin{block}{Exercise 05}
    Use mathematical induction, find the general term formula for the Fibonacci sequence defined recursively by:
    \[
    F_0 = 0, \quad F_1 = 1, \quad F_n = F_{n-1} + F_{n-2} \text{ for } n \ge 2.
    \]
\end{block}
\end{frame}

\section{Rational Numbers}
\begin{frame}{Rational Numbers and Basic Properties}
    The set of rational numbers is introduced from set theory as
    \[
    \mathbb{Q} := \left\{ \frac{p}{q} \;\middle|\; p,q \in \mathbb{Z}, q \ne 0 \right\}.
    \]
    Recap the basic properties of rational numbers:
    \vspace{0.2cm}
    \begin{columns}[t]
    \column{0.5\textwidth}
    \textbf{Properties for Addition:}
    \begin{itemize}
        \item Commutativity: $a + b = b + a$
        \item Associativity: $(a + b) + c = a + (b + c)$
        \item Existence of additive identity: $a + 0 = a$
        \item Existence of additive inverse: $a + (-a) = 0$
    \end{itemize}
    \column{0.5\textwidth}
    \textbf{Properties for Multiplication:}
    \begin{itemize}
        \item Commutativity: $a \cdot b = b \cdot a$
        \item Associativity: $(a \cdot b) \cdot c = a \cdot (b \cdot c)$
        \item Existence of multiplicative identity: $a \cdot 1 = a$
        \item Existence of multiplicative inverse: $a \cdot a^{-1} = 1$ for $a \ne 0$
        \item Distributivity: $a \cdot (b + c) = a \cdot b + a \cdot c$
    \end{itemize}
    \end{columns}
\end{frame}
\begin{frame}{Absolute Value and Triangle Inequality}
    \begin{definition}[Absolute Value]
        The \textbf{absolute value} (also called the \textbf{modulus}) of a rational number $x$ is defined as
        \[
        |x| := \begin{cases}
            x, & \text{if } x \ge 0, \\
            -x, & \text{if } x < 0.
        \end{cases}
        \]
    \end{definition}
    \begin{theorem}[Triangle Inequality]
        For any $x,y \in \mathbb{Q}$, the following inequality holds:
        \[
        |x + y| \le |x| + |y|.
        \]
        Inverse Triangle Inequality:
        \[
        ||x| - |y|| \le |x - y|.
        \]
    \end{theorem}
\end{frame}
\begin{frame}{Trichotomy Law and Closedness}
    We assume that we know what a \textbf{strictly positive} rational number is, i.e., that there exists a set $P$ with the property that for every number $a \in \mathbb{Q}$, one and only one of the following holds:
    \begin{enumerate}
        \item $a \in P$ (i.e., $a$ is strictly positive),
        \item $-a \in P$ (i.e., $a$ is strictly negative),
        \item $a = 0$.
    \end{enumerate}
    This is known as the \textit{\textbf{trichotomy law}} for rational numbers \textbf{(P10)}. It allows us to compare any two rational numbers and determine their order. \\
    \vspace{0.5cm}
    We also assume that the set of positive numbers is \textbf{closed} under addition and multiplication:
    \begin{enumerate}
        \item If $a,b \in P$, then $a + b \in P$. \textbf{(P11)}
        \item If $a,b \in P$, then $a \cdot b \in P$. \textbf{(P12)}
    \end{enumerate}
\end{frame}
%--------------------------------
\section{Real Numbers}

\begin{frame}{Bounds and Boundedness}

\begin{definition}[Upper and Lower Bounds]
Let $(X,\le)$ be a partially ordered set and $A \subseteq X$ be nonempty.  
\begin{itemize}
    \item An element $s \in X$ is an \textbf{upper bound} of $A$ if $a \le s$ for all $a \in A$.
    \item An element $s \in X$ is a \textbf{lower bound} of $A$ if $a \ge s$ for all $a \in A$.
\end{itemize}
\end{definition}

\begin{definition}[Bounded Sets in $\mathbb{Q}$]
A set \( U \subset \mathbb{Q} \) is said to be \textbf{bounded} if there exists a constant \( c \in \mathbb{Q} \) such that
\[
|x| \le c \text{ for all } x \in U.
\]
If this is not the case, we say that \( U \) is \textbf{unbounded}.
\end{definition}
We say that a set \( U \subset \mathbb{Q} \) is \textbf{bounded above} if there exists an upper bound
for \( U \), and \textbf{bounded below} if there exists a lower bound for \( U \).

\end{frame}

\begin{frame}{Maximum and Minimum, Supremum and Infimum}

\begin{definition}[Maximum and Minimum]
Let $U \subset \mathbb{Q}$.  
\begin{itemize}
    \item $x_1 \in U$ is the \textbf{minimum} of $U$ ($x_1 = \min U$), if $x_1 \le x$ for all $x \in U$.
    \item $x_2 \in U$ is the \textbf{maximum} of $U$ ($x_2 = \max U$), if $x_2 \ge x$ for all $x \in U$.
\end{itemize}
\end{definition}
\textbf{Note:} A set not bounded below has no minimum, and a set not bounded above has no maximum. \underline{Even a bounded set may not have a max or min.}
\begin{definition}[Supremum and Infimum]
Let $U \subset \mathbb{Q}$.  
\begin{itemize}
    \item An upper bound $c_2 \in \mathbb{Q}$ of $U$ is the \textbf{least upper bound} or \textbf{supremum}, denoted $c_2 = \sup U$, if no $c \in \mathbb{Q}$ with $c < c_2$ is an upper bound.
    \item A lower bound $c_1 \in \mathbb{Q}$ of $U$ is the \textbf{greatest lower bound} or \textbf{infimum}, denoted $c_1 = \inf U$, if no $c \in \mathbb{Q}$ with $c > c_1$ is a lower bound.
\end{itemize}
\end{definition}

\end{frame}


\begin{frame}{Greatest Lower and Least Upper Bounds of Sets}

\begin{definition}

We say that an upper bound \( c_2 \in \mathbb{Q} \) of a set \( U \subset \mathbb{Q} \) is
the \textbf{least upper bound} or \textbf{supremum} of \( U \) if no \( c \in \mathbb{Q} \) with \( c < c_2 \) is an
upper bound. We then write \( c_2 =: \sup U \).

\par\vspace{0.5em}
We say that a lower bound \( c_1 \in \mathbb{Q} \) is the \textbf{greatest lower bound} or
\textbf{infimum} of \( U \) if no \( c \in \mathbb{Q} \) with \( c > c_1 \) is a lower bound.
We then write \( c_1 =: \inf U \).
\end{definition}

\end{frame}


\begin{frame}{Naive Real Number}
\textbf{Problem: \underline{Rational numbers are insufficient.}}\\
\vspace{0.2cm}
\begin{centering}
    $\downarrow$ \\
\end{centering}
\vspace{0.2cm}
We define the set of real numbers \( \mathbb{R} \) as the smallest extension of the
rational numbers \( \mathbb{Q} \) such that the following property holds:

\begin{block}{Least Upper Bound Property (P13)}
If \( A \subset \mathbb{R} \), \( A \neq \emptyset \), and \( A \) is bounded above,  
then there exists a least upper bound for \( A \) in \( \mathbb{R} \).
\end{block}

We call all real numbers that are not rational ``\textbf{irrational numbers}''.
\end{frame}


\begin{frame}{Subsets of the Real Numbers}

\begin{definition}
Let \( A \subset \mathbb{R} \) be any set.

\begin{enumerate}[(i)]
\item We call \( x \in \mathbb{R} \) an \textbf{interior point} of \( A \) if there exists some \( \varepsilon > 0 \)
such that the interval \( (x - \varepsilon,\, x + \varepsilon) \subset A \).
The set of all interior points of \( A \) is denoted by \( \mathrm{int}\,A \).

\item We call \( x \in \mathbb{R} \) an \textbf{exterior point} of \( A \) if there exists some \( \varepsilon > 0 \)
such that \( (x - \varepsilon,\, x + \varepsilon) \cap A = \emptyset \).

\item We call \( x \in \mathbb{R} \) a \textbf{boundary point} of \( A \) if for every \( \varepsilon > 0 \),
both \( (x - \varepsilon,\, x + \varepsilon) \cap A \neq \emptyset \) and
\( (x - \varepsilon,\, x + \varepsilon) \cap A^{c} \neq \emptyset \).
The set of all boundary points of \( A \) is denoted by \( \partial A \).

\item We call \( x \in \mathbb{R} \) an \textbf{accumulation point} of \( A \) if for every \( \varepsilon > 0 \),
\( (x - \varepsilon,\, x + \varepsilon) \cap (A \setminus \{x\}) \neq \emptyset \).
\end{enumerate}
\end{definition}

\end{frame}


\begin{frame}{Subsets of the Real Numbers}

\begin{definition}
\begin{enumerate}
\item A set \( A \subset \mathbb{R} \) is called \textbf{open} if all points of \( A \) are interior points;  
that is, if \( A = \mathrm{int}\,A \).

\item A set \( A \subset \mathbb{R} \) is called \textbf{closed} if its complement \( \mathbb{R} \setminus A \) is open.

\item The set
\[
\overline{A} := A \cup \partial A
\]
is called the \textbf{closure} of \( A \).  
It can be shown that \( \overline{A} \) is the smallest closed set that contains \( A \).
\end{enumerate}
\end{definition}

\begin{block}{Exercise 06}
    Show that \(\mathbb{R} = \mathrm{int}\,A \;\cup\; \mathrm{ext}\,A \;\cup\; \partial A\), and think about the sets that are both open and closed.
\end{block}
\end{frame}


\begin{frame}{Subsets of the Real Numbers}

\begin{block}{Exercise 07}
    Proof that:
\begin{enumerate}[(a)]
\item For any collection \( \{ G_\alpha \} \) of open sets, the union \( \bigcup_\alpha G_\alpha \) is open.

\item For any collection \( \{ F_\alpha \} \) of closed sets, the intersection \( \bigcap_\alpha F_\alpha \) is closed.
\hfill (2.21)

\item For any finite collection \( G_1, G_2, \dots, G_n \) of open sets, the intersection
\(\bigcap_{i=1}^n G_i\) is open.

\item For any finite collection \( F_1, F_2, \dots, F_n \) of closed sets, the union
\(\bigcup_{i=1}^n F_i\) is closed.
\end{enumerate}
\end{block}
\end{frame}
\begin{comment}
\section{Complex Numbers}
\begin{frame}{Complex Numbers}
We define addition in \(\mathbb{C}\) by
\[
(a_1, b_1) + (a_2, b_2) := (a_1 + a_2,\, b_1 + b_2),
\]
so in particular the sum of two real numbers is again a real number.  
It is straightforward to verify that properties (P1)–(P4) hold for this addition.

Multiplication is defined as
\[
(a_1, b_1) \cdot (a_2, b_2) := (a_1 a_2 - b_1 b_2,\, a_1 b_2 + a_2 b_1).
\]
This multiplication satisfies axioms (P5)–(P8), and together with addition also satisfies (P9).  
The neutral element is \((1,0)\), as expected, while the inverse element is slightly more involved.

\end{frame}


\begin{frame}{Complex Numbers}

\textbf{Recall:} A binary relation \(R\) on a set \(S\) is a \textbf{total order} if it is
\begin{itemize}
    \item \textbf{Reflexive:} \(x R x\) for all \(x \in S\)
    \item \textbf{Anti-symmetric:} \(x R y\) and \(y R x \implies x = y\)
    \item \textbf{Transitive:} \(x R y\) and \(y R z \implies x R z\)
    \item \textbf{Connected:} For any \(x, y \in S\), either \(x R y\) or \(y R x\)
\end{itemize}

\textbf{Question:} Does there exist a total order on \(\mathbb{C}\)?  

\textbf{Observation:} People often say \emph{complex numbers cannot be ordered}.  

\textbf{Reason: Algebraic Compatibility}
\begin{itemize}
    \item Suppose there exists a total order \(\le\) on \(\mathbb{C}\) compatible with addition and multiplication.
    \item Consider the imaginary unit \(i\):
    \begin{itemize}
        \item If \(0 \le i\), then \(0 \le i^2 = -1\) — contradiction.
        \item If \(i \le 0\), then \(i^2 = -1 \ge 0\) — also contradiction.
    \end{itemize}
    \item Hence, no total order on \(\mathbb{C}\) can respect its algebraic structure.
\end{itemize}

\end{frame}

\begin{frame}{Complex Numbers}

\begin{definition}
The \textbf{modulus} or \textbf{absolute value} of \(z\) is defined by
\[
|z| := \sqrt{a^2 + b^2} = \sqrt{z \cdot \overline{z}},
\] 
\end{definition}

\begin{definition}
The \textbf{open ball} of radius \(R\) centered at \(z_0\) is
\[
B_R(z_0) := \{ z \in \mathbb{C} : |z - z_0| < R \}.
\]
A set \(\mathcal{S} \subset \mathbb{C}\) is called \textbf{bounded} if there exists some \(R > 0\) such that
\[
\mathcal{S} \subset B_R(0).
\]
\end{definition}

\end{frame}
\end{comment}

\begin{frame}{References and What's Next}
\large \textbf{References:}
\vspace{0.5em}
\begin{itemize}
  \item Horst Hohberger, \textit{math186\_all\_lecture\_slides (2026)}
  \item Wenxin Chen, \textit{MATH1860J RC}, 25 Fall
  \item Leqi Tang, \textit{MATH1860J RC}, 25 Fall
\end{itemize}
\vspace{1cm}
\large \textbf{To be continued...} \\
\vspace{0.5em}
\begin{itemize}
    \item Complex Numbers
    \item Functions and Maps
    \item Sequences
\end{itemize}
\end{frame}

\begin{frame}
    \begin{center}
    \huge \textbf{Thanks for Listening!} \\
    \vspace{0.5cm}
    \large \textbf{Any questions? Please feel free to ask!} \\
    \end{center}
\end{frame}

\end{document}
